\section{Conclusion}\label{sec:conclusion}
The central finding of this review is that the quantum-implicit tier is definable, its entry criteria are derivable from the existing literature and no current quantum programming system occupies it.
The tier hierarchy is populated at Assembly, IR, Framework and High-level, but the tier at which programmers write non-trivial programs without quantum-mechanical prerequisites remains absent.
This finding answers both research questions.
No existing system meets the requirements for the quantum-implicit tier (RQ1); the gap is structural rather than incremental and a new language targeting this tier is justified (RQ2).
The five contributions are a criterion-based tier taxonomy, a five-axis evaluation rubric with explicit score-4 entry conditions, an operationalisation of quantum bias, comparative assessments of eleven systems and empirical evidence that the quantum-implicit tier is unoccupied.

\paragraph{Limitations}
The evaluation rubric is theoretically grounded but has not been subjected to empirical validation.
Axis weights and score thresholds are derived from the literature and should be treated as principled rather than empirically calibrated~\cite{Coblenz2021PLIERS,Myers2016PLUsabilitySIG}.
The findings reflect a snapshot of the field as of 2026-04-08; the quantum programming ecosystem evolves rapidly and systems may advance beyond their assessed state by the time of publication.
The assessment is against theoretical abstraction requirements; practitioner evidence cited in Section~\ref{sec:background} is observational~\cite{ferreira2024qplusage} and a complementary user study would strengthen the design implications independently of the rubric.
The system selection criteria were pre-specified and PRISMA-compliant, but photonic computing languages and topological qubit frameworks are underrepresented; findings may not generalise fully across all quantum computing paradigms.

\paragraph{Rubric validation}
The most immediate future work is empirical validation of the evaluation instrument.
A documentation-only learnability study, in which participants complete T1--T5 using only public documentation for a small set of systems, would provide observable proxies for quantum bias burden~\cite{Coblenz2021PLIERS}.
An inter-rater reliability study establishing Cohen's $\kappa$ across independent reviewers would provide a formal reliability measure for the axis scores~\cite{Myers2016PLUsabilitySIG}.

\paragraph{System set and task suite expansion}
Systems excluded for failing a single inclusion criterion and systems that post-date the protocol registration are natural candidates for an expanded review.
Domain-specific task suites covering quantum chemistry, variational algorithms and quantum networking would stress-test the rubric beyond the general T1--T5 workloads~\cite{ferreira2024qplusage}.
The Assembly-tier boundary is also evolving; OpenQASM's progression toward intermediate representation features suggests that tier boundaries themselves may require revision as the field matures.

\paragraph{Formalisation directions}
Formalising the semantics of quantum-implicit constructs is the theoretical counterpart to this empirical study.
The foundational abstractions framework~\cite{yuan2025foundationalabstractions} provides the most direct starting point for specifying compiler obligations at the quantum-implicit tier.
Verified compilation patterns for quantum domains demonstrate that end-to-end correctness guarantees across a compilation pipeline are achievable~\cite{verifiedcompiler2025qblue}.
ZX-calculus completeness results establish a graphical framework capable of verifying any quantum circuit equivalence~\cite{zxworkingquantumscientist2020}, providing a candidate theoretical foundation for semantics-preserving compilation rewrites.

\paragraph{Prototype language design}
The design implications in Section~\ref{subsec:discussion:implications} motivate a prototype language that embodies score-4 on all five axes while retaining explicit escape hatches for optimisation and hardware-specific tuning.
Designing, implementing and evaluating such a language against the T1--T5 task suite is the immediate next step in the research programme to which this review contributes.